% =====================================================================
% 13_polarization_observables.tex
% 第 13 章 极化观测量的部分波合成
% Wave 4 / Agent (ch13)
% =====================================================================

\chapter{极化观测量的部分波合成}
\label{ch:13_polarization_observables}

% ============================================================
\section{物理动机}
\label{sec:ch13-motivation}

\paragraph{为什么不止看截面？}
第 \ref{ch:12_observables_basics} 章把 AGS 求出的 \(\Uop\)-振幅
\(\langle\PWchanp|\Uop|\PWchan\rangle\) 翻译成了
\emph{自旋平均的}（spin-averaged）弹性微分截面 \(\diffXS(\theta)\)。
对一个标量靶（如 \(p+p\)、\(\alpha+\alpha\)）这就是故事的全部；但对
\(d+p\) 这种「自旋-1 + 自旋-1/2 → 自旋-1 + 自旋-1/2」的散射，自旋自由度
还能携带远比截面更精细的信息。原因有两层：

\begin{enumerate}
  \item \textbf{张量力是氘核存在的根本原因}。一个仅含中心力（\(\Vop \propto
        \mathbb{1}\)）的核子-核子势永远束缚不出 \(I=0\)、\(J=1\) 的氘核——
        因为 \({}^3S_1\) 单独的吸引深度不够，需要张量算符
        \(S_{12}(\hat{r}) \equiv 3(\bm{\sigma}_1\cdot\hat{r})(\bm{\sigma}_2\cdot\hat{r})
        - \bm{\sigma}_1\cdot\bm{\sigma}_2\) 把 \({}^3S_1\) 与
        \({}^3D_1\) 耦合起来，靠 \(D\) 态混入贡献的 \(\sim 4\%\) 把束缚拉
        到 \(\Ed = -2.2245\,\text{MeV}\)。同样的张量结构在 \(d+p\) 散射中产
        生 \emph{自旋上的角度依赖}：散射后氘核会有不同方向的
        \(M_J\) 占据率，反映在张量分析能力 \(\Ttwentyzero, \Ttwentyone,
        \Ttwentytwo\)。
  \item \textbf{三核子力（3NF）首先在极化里露脸}。Hannover/Krakow 系列工作
        \cite{Witala2008} 反复显示一个事实：在 \(\Tlab \gtrsim 100\,\text{MeV}/u\) 时，
        二体力（NN-only）对 \(\diffXS\) 的预言已经能与实验对到 5\%，但对
        \(A_y\)（质子矢量分析能力）和 \(\iTeleven\)（氘核矢量分析能力）会出现
        系统性的「\(A_y\) puzzle」——理论预言的极值大约只有实验的 70\%。
        加进 \(\Delta\)-isobar / 2\(\pi\)-exchange 3NF 后差距收敛。这意味
        着 \emph{极化观测量是 3NF 物理学的第一探针}。
\end{enumerate}

第 \ref{ch:15_3nf_physics} 章会专门讲三核子力的物理与部分波投影，本章只把
极化观测量从 \(\Uop\)-振幅推到「角度曲线」这一段管道铺好。

\paragraph{从 \(M\) 到 \(T_{kq}\)：自旋张量分解的视角。}
对一个一般的散射过程
\begin{equation}
\label{eq:ch13-scattering-process}
\ket{a, M_a; b, M_b}_{\text{in}} \;\longrightarrow\;
\ket{a', M_{a'}; b', M_{b'}}_{\text{out}},
\end{equation}
振幅矩阵
\(M_{M_{a'}M_{b'}; M_a M_b}(\theta,\phi)
\equiv \langle a' M_{a'}, b' M_{b'}|\, T(\theta,\phi)\,|a M_a, b M_b\rangle\)
是一个张量积矩阵。当我们把入射粒子的极化态用密度矩阵
\(\rho_{a}^{\text{in}}\) 描述、出射做无极化求和时，微分截面写为迹
\begin{equation}
\label{eq:ch13-diff-trace}
\diffXS(\theta,\phi)
\;=\; \frac{1}{2 s_a + 1}\,\frac{1}{2 s_b + 1}\,
\Tr\!\left[ M(\theta,\phi)\, \rho^{\text{in}}\, M^\dagger(\theta,\phi)\right],
\end{equation}
其中 \(\rho^{\text{in}} = \rho_a^{\text{in}} \otimes \rho_b^{\text{in}}\)。
若入射 \(a\)（氘核）有矢量极化 \(p_y \equiv P_y\) 与张量极化 \(p_{kq}\)
（\(k=2\)），则
\begin{equation}
\label{eq:ch13-rho-deuteron}
\rho_a^{\text{in}}(\bm{p}_a)
\;=\;
\frac{1}{3}\sum_{k=0}^{2}\sum_{q=-k}^{k}
(2 k + 1)\, p_{kq}^{*}\, \tau_{kq},
\end{equation}
而 \(\tau_{kq}\) 是\textbf{自旋张量算符}（spin tensor operators）
\begin{equation}
\label{eq:ch13-tau-def}
\langle s_a M' | \tau_{kq} | s_a M\rangle
\;=\; \sqrt{2 k + 1}\,(-1)^{s_a - M'}\,
\begin{pmatrix} s_a & k & s_a \\ -M' & q & M \end{pmatrix}
\;\equiv\; (-1)^{s_a - M'}\,
\langle s_a, -M' ; s_a, M | k, q\rangle,
\end{equation}
满足正交关系
\(\Tr[\tau_{kq}^\dagger \tau_{k'q'}] = (2 s_a + 1)\,\delta_{kk'}\delta_{qq'}\)。
这套基张开了氘核 \(3 \times 3\) 自旋空间的所有 Hermitian 算符
（\(1 + 3 + 5 = 9 = 3^2\) 个分量）。本章下面 \cref{sec:ch13-math} 会把式
\eqref{eq:ch13-diff-trace} 拆成「极化无关 \(+\) 矢量分析能力 \(+\) 张量分
析能力」三块。

\paragraph{本章在 Tic-tac 流水线里的位置。}
\(\Uop\)-矩阵元由第 \ref{ch:11_faddeev_solver} 章 Padé 加速 Neumann 迭
代输出；第 \ref{ch:12_observables_basics} 章把它合成成弹性振幅
\(M_{M_{a'}M_{b'}; M_a M_b}\) 的部分波分解；本章把这个 \(M\)-矩阵进一步
合成成 \(\diffXS\)、\(\Ay\)、\(\iTeleven\)、\(\Ttwentyzero\)、
\(\Ttwentyone\)、\(\Ttwentytwo\) 五条角度曲线，并提供与实验比较的代码骨
架（\(190\,\text{MeV}\) 基准）。第 \ref{ch:14_miller_gate1} 章会用相同
的 \(M\)-矩阵推 Miller-gate 1 相移；第 \ref{ch:15_3nf_physics} 章及之后
通过\textbf{切换 3NF 开关}对比本章的曲线，定量回答「3NF 改进了多少
\(\iTeleven\)？」。

\paragraph{读完本章你应该能。}
\begin{itemize}
  \item 在脑中默写式 \eqref{eq:ch13-rho-deuteron} 与张量极化的 5 个分量
        \(\Ttwentyzero, \Re\Ttwentyone, \Im\Ttwentyone, \Re\Ttwentytwo,
        \Im\Ttwentytwo\) 之间的关系，并解释为何只有「实部 \(\Re\Ttwentyone\)
        与 \(\Im\Ttwentytwo\)」在 Madison 约定下出现在散射平面上。
  \item 准确写出 \(\Ay^{(p)}\)（质子作为入射极化粒子）与 \(\iTeleven\)
        （氘核作为入射极化粒子）之间的因子 \(\sqrt{3/2}\) 关系，并知道在
        Sekiguchi et al.\,的 \(190\,\text{MeV}\) 数据表里写的是哪一个量。
  \item 拿到一个 \texttt{U\_PW\_elements\_*.txt} 文件后，能用
        \texttt{compare\_Ay\_experiment.py} 的 \texttt{parse\_u\_file} +
        \texttt{combine\_channels\_by\_energy} +
        \texttt{\_predict\_observables\_from\_u} 三步把它还原为五条曲线。
  \item 解释当前实现是「确定性 reduced-U 映射」（deterministic reduced-U
        map），与教科书里「自旋张量逐分量求迹」（full spin tensor trace）
        在数值上相差什么、为什么这样写还能在 \(190\,\text{MeV}\) 基准里达
        到 \(\iTeleven\,\text{RMSE} \le 0.10\) 的早期通过线。
\end{itemize}

% ============================================================
\section{数学推导}
\label{sec:ch13-math}

\subsection{球张量基与极化态密度矩阵}
\label{subsec:ch13-spherical-tensor}

球张量算符 \(\tau_{kq}\)（\(k=0,\ldots,2 s_a;\;q=-k,\ldots,k\)）在自旋
\(s_a\) 子空间的完备性使任意 Hermitian \(\rho_a\) 都能用
\(\rho_a = \frac{1}{2 s_a + 1}\sum_{kq}(2 k + 1)\,p_{kq}^{*}\,\tau_{kq}\)
来展开。\(p_{kq}\) 通常被分别命名为：
\begin{itemize}
  \item \(k=0\)：\(p_{00} = 1\)（密度矩阵迹归一）
  \item \(k=1\)：\textbf{矢量极化} \(p_{1\pm 1}, p_{10}\)，物理意义是
        \(\langle s_y\rangle\) 等
  \item \(k=2\)：\textbf{张量极化} \(p_{2 q}\)，对应 \(\langle s_i s_j\rangle\)
        的对称无迹组合
\end{itemize}
对氘核 \(s_a = 1\)，张量分量数为 \(2k+1=5\)。

\begin{remark}
\label{rem:ch13-pkq-vs-Tkq}
约定上 \(p_{kq}\) 是入射粒子\textbf{极化度}（极化态的几何描述），\(T_{kq}\)
是同名球张量基下散射后\textbf{出射粒子的极化矩}（散射动力学输出）。两者在
公式形状上对称，但物理身份不同：
\(p_{kq}\) 由束流极化设备决定（人为设定），\(T_{kq}\) 由 \(\Uop\)-振幅决定
（理论预言）。\textbf{分析能力 \(\iTeleven\)、\(\Ttwentyzero\) 等是「以入
射极化态 \(p_{kq}\) 为探针、出射对 \(p_{kq}\) 的响应灵敏度」}，因此名字带
着 \(T\) 是暗示它们是\emph{响应函数}，而非真实测得的张量极化矩本身。
\end{remark}

\subsection{从 \(M\)-矩阵到极化观测量的求迹公式}
\label{subsec:ch13-M-to-Tkq}

把 \eqref{eq:ch13-rho-deuteron} 代入 \eqref{eq:ch13-diff-trace}，做一些
直接的张量运算可得
\begin{equation}
\label{eq:ch13-master-trace}
\diffXS(\theta,\phi) \cdot \big[1 + p_{kq}\, T_{kq}^{(s_a)}(\theta) + \cdots\big]
\;=\;
\frac{1}{(2 s_a + 1)(2 s_b + 1)}
\Tr\!\left[ M\, \rho^{\text{in}}\, M^\dagger\right],
\end{equation}
其中
\begin{equation}
\label{eq:ch13-Tkq-master}
\boxed{\;
T_{kq}^{(s_a)}(\theta)
\;=\;
\frac{\Tr\!\left[ M\, \tau_{kq}^{(a)}\,M^\dagger\right]}
     {\Tr\!\left[ M\, M^\dagger\right]}.
\;}
\end{equation}
这是教科书 \cite{OhlsenReports} 形式上最干净的「分析能力 = 对应自旋张量
算符在散射前后的 \(M\)-矩阵相干平均」。求迹是\textbf{在自旋指标 \(M_a, M_b,
M_{a'}, M_{b'}\) 上} 做的；散射角依赖来自 \(M(\theta,\phi)\) 自身。

\paragraph{矢量分析能力 \(\iTeleven\)（\(k = 1, q = +1\)）。}
\begin{equation}
\label{eq:ch13-iT11-formula}
\iTeleven(\theta)
\;=\;
\frac{i\,\Tr\!\left[ M\, \tau_{1,+1}^{(a)}\,M^\dagger\right]}
     {\Tr\!\left[ M\, M^\dagger\right]},
\qquad
\tau_{1,+1}^{(a)} \;=\; -\frac{1}{\sqrt{2}}(s_x^{(a)} + i\, s_y^{(a)}).
\end{equation}
这是\textbf{氘核}（\(s_a = 1\)）的矢量分析能力；前置因子 \(i\) 来自约定，
保证 \(\iTeleven\) 是实数。它在 Madison 约定下与「氘核入射时沿
\(\hat{n}\equiv \hat{k}_{\text{in}}\times\hat{k}_{\text{out}}/\|\cdots\|\)
方向纵向极化、再测量散射偏左/偏右不对称」给出的实验量相同。

\paragraph{张量分析能力 \(\Ttwentyzero, \Ttwentyone, \Ttwentytwo\)。}
\begin{equation}
\label{eq:ch13-T20-formula}
\Ttwentyzero(\theta)
\;=\;
\frac{\Tr\!\left[ M\,\tau_{2,0}^{(a)}\,M^\dagger\right]}
     {\Tr\!\left[ M\,M^\dagger\right]},
\qquad
\tau_{2,0}^{(a)} = \frac{1}{\sqrt{6}}\left(3 (s_z^{(a)})^2 - s_a(s_a+1)\,\mathbb{1}\right).
\end{equation}
\begin{equation}
\label{eq:ch13-T2pm1-formula}
\Ttwentyone(\theta)
\;=\;
\frac{\Tr\!\left[ M\,\tau_{2,+1}^{(a)}\,M^\dagger\right]}
     {\Tr\!\left[ M\,M^\dagger\right]},
\qquad
\tau_{2,+1}^{(a)} = -\tfrac{1}{2}\left(s_z^{(a)} s_+^{(a)} + s_+^{(a)} s_z^{(a)}\right),
\end{equation}
\begin{equation}
\label{eq:ch13-T22-formula}
\Ttwentytwo(\theta)
\;=\;
\frac{\Tr\!\left[ M\,\tau_{2,+2}^{(a)}\,M^\dagger\right]}
     {\Tr\!\left[ M\,M^\dagger\right]},
\qquad
\tau_{2,+2}^{(a)} = \tfrac{1}{2}(s_+^{(a)})^2.
\end{equation}
其中 \(s_\pm = s_x \pm i s_y\)。Madison 约定（见
\cref{subsec:ch13-Madison-vs-PRC}）下散射平面取 \(xz\)，故
\(\Ttwentyone\) 与 \(\Ttwentytwo\) 的「\(q\) 取虚部还是实部」是带物理意
义的：实验报告的 \(\Ttwentyone\) 是 \(\Re\Ttwentyone\)，
\(\Ttwentytwo\) 是 \(\Re\Ttwentytwo\)（两者在实验上都是实数）。

\subsection{极化无关分量与 \(M\)-矩阵的简化记法}
\label{subsec:ch13-helicity-amplitudes}

对 \(d + p\)（\(s_a = 1, s_b = 1/2\)），\(M\)-矩阵自旋空间维数是
\((2 s_a+1)(2 s_b+1) = 6\)，所以 \(M\) 是 \(6\times 6\) 复矩阵。Faddeev
求解器输出的「\(\Uop\)-amplitude」沿\textbf{部分波展开}写下，每个
\((J^\pi, T)\) 块给一个 \(2\times 2\)（\(j_{2N} = 1\) 通道：\({}^3S_1\)
与 \({}^3D_1\) 的 D-S 耦合）矩阵 \(\Uop_{\PWchanp\PWchan}^{J^\pi T}\)。
合成回 \(M_{M_{a'}M_{b'}; M_a M_b}(\theta)\) 时要乘以 Wigner-D 函数
与 Clebsch-Gordan 系数；这一步在
\cref{ch:12_observables_basics} 已经详细推导。

\textbf{当前 Tic-tac 实现的简化处理}：在 \texttt{compare\_Ay\_experiment.py}
里，每个能量给定时只取
\((J^\pi = 1^+, 1^-)\) 两个通道的「\(2\times 2\) 缩并 U-矩阵」
\(U_{ij} \equiv U^{J^\pi}_{ij}\)（\(i, j \in \{0, 1\}\) 对应
\({}^3S_1, {}^3D_1\) 的归一化基），按权 \(\sigma_{\text{proxy}}\) 平均后构
造\textbf{reduced-U 不变量}：
\begin{align}
\label{eq:ch13-reduced-U-norm}
u_{\text{norm}} &= \sum_{ij} |U_{ij}|^2, \\
\label{eq:ch13-inv-iT11}
\text{inv\_iT11} &= \frac{\Im\!\left[(U_{00} + U_{11})^{*}(U_{01} - U_{10})\right]}{u_{\text{norm}}}, \\
\label{eq:ch13-inv-T20}
\text{inv\_T20} &= \frac{|U_{00}|^2 + |U_{11}|^2 - |U_{01}|^2 - |U_{10}|^2}{u_{\text{norm}}}, \\
\label{eq:ch13-inv-T21}
\text{inv\_T21} &= \frac{\Re\!\left[(U_{00} - U_{11})^{*}(U_{01} + U_{10})\right]}{u_{\text{norm}}}, \\
\label{eq:ch13-inv-T22}
\text{inv\_T22} &= \frac{\Re\!\left[U_{00}^{*}\,U_{11} - U_{01}^{*}\,U_{10}\right]}{u_{\text{norm}}}.
\end{align}
然后把它们投影到 Legendre 多项式 \(P_n(\cos\theta)\) 上得角度曲线（
\cref{eq:ch13-iT11-pred,eq:ch13-dsigma-pred}）。这是\emph{显式可解析的、
不需要 Wigner-D 求和的近似映射}，作用是把「U-矩阵元的相位/幅度信息」翻
译成「\(\theta\) 上的曲线形状」，同时保留 \(\iTeleven\) 反对称（与
\(\Im\) 项绑死）、\(\Ttwentyzero\) 对称（与对角范数差绑死）等基本结构。
完整教科书形式将在 \cref{ch:14_miller_gate1} 之后的 J-tensor 实现里上线。

\subsection{\(\Ay\) 与 \(\iTeleven\) 的因子 \(\sqrt{3/2}\)}
\label{subsec:ch13-Ay-vs-iT11}

实验上经常会同一篇论文里同时报告 \(\Ay\)（质子矢量分析能力）与
\(\iTeleven\)（氘核矢量分析能力）。两者来自完全相同的物理（自旋-轨道耦
合 + 张量力），但归一不同：
\begin{equation}
\label{eq:ch13-Ay-iT11}
\boxed{\;
\Ay^{(d)}(\theta) \;=\; \sqrt{\frac{3}{2}}\;\iTeleven(\theta).
\;}
\end{equation}
推导路径：氘核 \(s_a = 1\) 时 \(\tau_{1,+1}^{(a)} = -\frac{1}{\sqrt{2}}(s_x +
i s_y)\)，而 \(\Ay\) 的标准定义直接是 \(\langle s_y\rangle\) 的反对称
（无 \(\frac{1}{\sqrt{2}}\) 与 \(\sqrt{2k+1}\) 的归一），算下来差一个
\(\sqrt{3/2}\)。\textbf{Sekiguchi et al.\,（2002, PRC 65, 034003）报的
是 \(\iTeleven\)，而老一些的 RIKEN 数据表中有时直接写 \(\Ay\)；混淆这两
者会得到 22.5\% 的系统偏差}。第 \ref{subsec:ch13-pitfall-Ay-vs-iT11} 节
专门强调这一陷阱。

\subsection{Madison 约定 vs PRC 约定}
\label{subsec:ch13-Madison-vs-PRC}

\textbf{Madison 约定}（1971，Univ. of Wisconsin）规定：
\begin{itemize}
  \item \(\hat{z}\) 沿入射束方向 \(\hat{k}_{\text{in}}\)；
  \item \(\hat{y} = \hat{k}_{\text{in}}\times \hat{k}_{\text{out}} / |\cdots|\)
        垂直散射平面；
  \item \(\hat{x} = \hat{y}\times \hat{z}\) 闭合右手系；
  \item 球张量分量 \(\tau_{kq}\) 用 \(\hat{z}\) 量化。
\end{itemize}
PRC（\textit{Phys. Rev. C}）的现代约定与 Madison 一致，但部分老论文
（特别是 70 年代的 \textit{Phys. Lett. B}）把 \(\hat{y}\) 取为
\(\hat{k}_{\text{out}}\times \hat{k}_{\text{in}}\)（反向），这会让
\(\iTeleven\) 整体翻号。Tic-tac 的 \texttt{compare\_Ay\_experiment.py}
默认使用 Madison/PRC 约定（\(\iTeleven > 0\) 在前向中等角度），与
Sekiguchi et al.\ 表 III 一致。

% ============================================================
\section{离散化}
\label{sec:ch13-discretization}

第 \ref{ch:12_observables_basics} 章已经把
\(\Uop\)-振幅 \(\to\) \(M\)-矩阵的部分波合成的离散网格交代清楚：

\begin{itemize}
  \item \textbf{动量网格}：\(p\) (\({}^3\!S\) 系内部相对动量) 与 \(q\)
        （第三粒子相对于 pair 12 的动量）共用第 \ref{ch:06_wp_swp_states}
        章的 \(\NpWP \times \NqWP\) WP/SWP 网格；
  \item \textbf{角度网格}：\(\theta_{\text{cm}} \in [0°, 180°]\) 通常按
        实验数据点取（无须自定义网格），见
        \cref{lst:ch13-angles-from-exp}；
  \item \textbf{能量网格}：\(\Tlab\) 给定时只在「最近的 on-shell q-bin
        中点」上估算，详见
        \cref{ch:11_faddeev_solver} 与
        \cref{subsec:ch13-pitfall-input-vs-solver-tlab}。
\end{itemize}

本章额外引入的离散化只有一项——\textbf{角度积分网格}。当前实现里
\texttt{compare\_Ay\_experiment.py} 不做角度积分（它直接在实验给出的 26
个角度点上求 \(M(\theta_i)\)）。如果要算\textbf{角度积分截面}
\(\sigma_{\text{tot}}\)，应当用 Gauss-Legendre 积分 \(N_\theta = 24\)
点，足以收敛到 \(10^{-4}\)（\cref{app:C_quadrature} 附录 C）。

\begin{lstlisting}[style=python, caption={从 T.txt 取实验角度网格}, label={lst:ch13-angles-from-exp}]
# data/DataOfCrosssectionAndPol/CompletSetOFT/T.txt 的第一列就是
# theta_cm[deg]; compare_Ay_experiment.py 直接复用这套网格
exp_it11 = read_experimental_iT11(tensor_data)  # angles, values, errors
exp_dsigma = read_experimental_dsigma(dsigma_data, target_unit=...)
# 之后所有模型预言都在 exp_it11["angles"] 与 exp_dsigma["angles"] 上求值
\end{lstlisting}

% ============================================================
\section{算法骨架：POLARIZATION-OBSERVABLES}
\label{sec:ch13-algorithm}

\begin{algorithm}[H]
\caption{\textsc{Polarization-Observables}\((\Tlab^{\text{target}},
\text{run\_params})\)}
\label{alg:ch13-pol-obs}
\KwIn{目标实验室能量 \(\Tlab^{\text{target}}\)（MeV），网格参数
\((\NpWP, \NqWP, j_{2N}^{\max}, J_{3N}^{\max}, \ldots)\)}
\KwOut{曲线 \(\{\diffXS(\theta_i)\}, \{\iTeleven(\theta_i)\},
\{\Ttwentyzero(\theta_i)\}, \{\Ttwentyone(\theta_i)\},
\{\Ttwentytwo(\theta_i)\}\)；摘要 JSON / CSV / PNG}
\BlankLine
\textbf{Phase 1 — Solver 端} (\texttt{deuteron\_proton\_Ay.py})\;
\quad 1. 写出 solver 输入文件
\(\texttt{input\_tlab\_<label>\_solver\_validation\_autogen.txt}\)\;
\quad 2. 调用 \texttt{CPP/run input.txt}（产生 \(\Uop\)-振幅）\;
\quad 3. 校验 \texttt{U\_PW\_elements\_*.txt} 里有可解析的复数
        条目（\texttt{has\_u\_data}）\;
\BlankLine
\textbf{Phase 2 — 实验端} (\texttt{read\_experimental\_*})\;
\quad 4. 解析 \texttt{T.txt}：得 \(\big\{\theta_i^{(\text{exp})},\,
        (\iTeleven)_i^{(\text{exp})},\, (\Delta\iTeleven)_i\big\}\)\;
\quad 5. 解析 \texttt{DSigamaOverDOmega.txt}：得
        \(\big\{\theta_j^{(\text{ds})},\, (\diffXS)_j^{(\text{exp})}\big\}\)；按
        \texttt{--dsigma-unit} 转单位\;
\BlankLine
\textbf{Phase 3 — 合成端} (\texttt{compare\_Ay\_experiment.py})\;
\quad 6. \texttt{parse\_u\_file}：\texttt{U\_PW\_elements\_*.txt}
        \(\to [\text{SolverChannelPoint}]\)；按文件名解析
        \((J^\pi = 2J/\pm)\)\;
\quad 7. \texttt{combine\_channels\_by\_energy}：相同 \(\Tlab\) 上对宇
        称两通道做 \(\sigma_{\text{proxy}}\)-加权 \(\Uop\) 平均\;
\quad 8. 对每个能量 \(\Tlab_{(e)}\)，
        \texttt{\_predict\_observables\_from\_u}：\;
\qquad  8.1 计算 \(u_{\text{norm}}\)、\(\text{inv\_iT11}\)、
            \(\text{inv\_T20}\)、\(\text{inv\_T21}\)、\(\text{inv\_T22}\)
            （\cref{eq:ch13-reduced-U-norm,eq:ch13-inv-iT11,eq:ch13-inv-T20,eq:ch13-inv-T21,eq:ch13-inv-T22}）\;
\qquad  8.2 角度投影：
            \(\iTeleven^{\text{model}}(\theta) =
              \mathrm{clamp}(1.20\,\text{inv\_iT11}\,\sin\theta\,
              (-1.20\,P_2 - 0.20\,P_1), -1, 1)\)
            （\cref{eq:ch13-iT11-pred}）\;
\qquad  8.3 截面投影：
            \(\diffXS^{\text{model}}(\theta) = u_{\text{norm}}\,
              \exp(1.10\,\text{inv\_T20}\,P_2 +
              0.60\,\text{inv\_T22}\,P_4)\)
            （\cref{eq:ch13-dsigma-pred}）\;
\quad 9. 选离 \(\Tlab^{\text{target}}\) 最近的求解器能量为
        \texttt{best\_energy}\;
\quad 10. 写 \texttt{solver\_validation\_tlab\_<label>.\{json,txt\}}、
         \texttt{best\_energy\_*\_curve.csv}、SVG 与（plot\_validation\_curves
         可选）PNG\;
\end{algorithm}

\paragraph{阶段拆分背后的工程动机。}
把 Solver / 实验 / 合成三阶段拆开\textbf{是为了缓存复用}。Faddeev 求解
器对一个 \(\NpWP=20, \NqWP=20\) 网格在 4 线程下大约要 5 分钟（一次
\(P_{123}\) 计算 + Padé 迭代 12 步），而合成阶段在毫秒级；所以一旦
\texttt{U\_PW\_elements\_*.txt} 写好，研究者可以\emph{自由切换}
观测量公式（reduced-U → 完整张量求迹）、自由调
\texttt{--dsigma-unit}、自由换实验数据集，不用重跑求解器。这也是
\texttt{compare\_Ay\_experiment.py} 把 \texttt{--regenerate} 做成显式开关
的原因。

% ============================================================
\section{代码落地}
\label{sec:ch13-code}

\subsection{\texttt{deuteron\_proton\_Ay.py}：求解器驱动}
\label{subsec:ch13-code-driver}

\paragraph{Step A —— 写输入文件。}
\texttt{write\_solver\_inputs} 把 \texttt{SolverRunConfig} 数据类的字段拍平
成一份「自动生成的求解器配置文件」，其格式与第 \ref{ch:appE} 章附录 E 描
述的纯键-值表完全一致。注意能量单独写到一份
\texttt{tlab\_<label>\_validation\_autogen.txt}（求解器消费的 ASCII 列），
而 solver 入口取的是另一份配置文件指向它的 \texttt{energy\_input\_file}
路径——这个间接层让\textbf{多能量批处理}（如 70/135/190 三能量扫描）只需
改一个文件就行。

\lstinputlisting[style=python,
  caption={\texttt{deuteron\_proton\_Ay.py} 的 \texttt{write\_solver\_inputs}：
           生成 \(\Uop\)-计算所需的输入对},
  label={lst:ch13-write-solver-inputs}]
  {code_excerpts/snippets/ch13_write_solver_inputs.py}

\paragraph{Step B —— 调用求解器并校验输出。}
\texttt{run\_cpp\_solver} 把 \texttt{cwd} 切到 \texttt{CPP/}，让
\texttt{HDF5\_DISABLE\_VERSION\_CHECK=2} 屏蔽 conda 环境与系统
\texttt{libhdf5} 之间的小版本差异，把所有 stderr 与 stdout 都流到
\texttt{solver\_run.log}，最后扫描
\texttt{U\_PW\_elements\_*.txt} 检查是否含 7 列以上的可解析复数行。

\lstinputlisting[style=python,
  caption={\texttt{deuteron\_proton\_Ay.py} 的 \texttt{run\_cpp\_solver}：
           调度 \texttt{CPP/run} + 输出验证},
  label={lst:ch13-run-cpp-solver}]
  {code_excerpts/snippets/ch13_run_cpp_solver.py}

\subsection{\texttt{compare\_Ay\_experiment.py}：合成与比较}
\label{subsec:ch13-code-compare}

\paragraph{Step C —— 解析 \(\Uop\)-文件并构造 \(\sigma_{\text{proxy}}\)、
\(A_y^{\text{proxy}}\)。}
\texttt{parse\_u\_file} 一次性读取一份 \texttt{U\_PW\_elements\_*.txt}，
按文件名识别 \(J^\pi\)（\texttt{detect\_two\_j\_from\_filename} 解析
\verb|_JP_<2J>_<P>_|；\texttt{detect\_parity\_symbol} 取 \(\pm\)），
把每行的 \((U_{00}, U_{01}, U_{10}, U_{11})\) 复数四元组打包进
\texttt{SolverChannelPoint}。同时它在每行内部计算两个\textbf{诊断量}：
\begin{align}
\label{eq:ch13-fnoflip-fflip}
f_{\text{no-flip}}(p, q, e) &\equiv U_{00}(p, q, e) + U_{11}(p, q, e),\\
f_{\text{flip}}(p, q, e)    &\equiv U_{01}(p, q, e) + U_{10}(p, q, e),\\
\sigma_{\text{proxy}}(p, q, e) &\equiv |f_{\text{no-flip}}|^2 + |f_{\text{flip}}|^2,\\
A_y^{\text{proxy}}(p, q, e) &\equiv
\frac{\Im\![f_{\text{no-flip}}^{*} f_{\text{flip}}]}{\sigma_{\text{proxy}}}.
\end{align}
这两个 proxy 量\textbf{只用作通道加权与合并选择}，不用于最终的
\(\iTeleven, \Ttwentyzero\) 等观测量预言。

\lstinputlisting[style=python,
  caption={\texttt{compare\_Ay\_experiment.py} 的 \texttt{parse\_u\_file}：
           按 \(J^\pi\) 块解析 U-文件并构造 SolverChannelPoint},
  label={lst:ch13-parse-u-file}]
  {code_excerpts/snippets/ch13_parse_u_file.py}

\paragraph{Step D —— 实验数据读取。}
\texttt{read\_experimental\_iT11} 与 \texttt{read\_experimental\_dsigma}
都对原始 ASCII 表格做：跳过表头/注释、跳过 null 行、按列号解析。注意
\texttt{read\_experimental\_dsigma} 自动检测原始单位
（\texttt{mb/sr} 或 \texttt{fm2/sr}），并通过
\texttt{convert\_dsigma\_series} 转到 \texttt{--dsigma-unit} 指定的输出单位。

\lstinputlisting[style=python,
  caption={\texttt{compare\_Ay\_experiment.py} 的实验数据读取：iT11 与
           \(\diffXS\)},
  label={lst:ch13-read-experiment}]
  {code_excerpts/snippets/ch13_read_experiment.py}

\paragraph{Step E —— 通道合并。}
\texttt{combine\_channels\_by\_energy} 对相同 \(\Tlab\)（精度 \(10^{-6}\)
MeV）做 \(\sigma_{\text{proxy}}\) 加权平均，把
\((J^\pi = 1^+, J^\pi = 1^-)\) 这两条同能量、不同宇称的通道合成一组「等
效 U-矩阵」\((U_{00}^{\text{eff}}, U_{01}^{\text{eff}}, U_{10}^{\text{eff}},
U_{11}^{\text{eff}})\)。这一步是当前实现的「简化合成」，对应
\cref{subsec:ch13-helicity-amplitudes} 末段所说的「跳过 Wigner-D 求和的
近似映射」；完整实现里这里应替换成「对所有 \((J^\pi, T)\) 块求 Wigner-D
缩并」。

\lstinputlisting[style=python,
  caption={\texttt{compare\_Ay\_experiment.py} 的
           \texttt{combine\_channels\_by\_energy}：能量分组 +
           \(\sigma_{\text{proxy}}\)-加权 \(\Uop\) 平均},
  label={lst:ch13-combine-channels}]
  {code_excerpts/snippets/ch13_combine_channels.py}

\paragraph{Step F —— 不变量 + 角度投影。}
\texttt{\_predict\_observables\_from\_u} 是数学公式
\eqref{eq:ch13-reduced-U-norm}\nobreakdash--\eqref{eq:ch13-inv-T22} 的直接
代码版本：先算 4 个 reduced-U 不变量；再分别投到
\begin{align}
\label{eq:ch13-iT11-pred}
\iTeleven^{\text{model}}(\theta) &=
\mathrm{clamp}\!\left[\,1.20\,\text{inv\_iT11}\,\sin\theta\,
\big(-1.20\,P_2(\cos\theta) - 0.20\,P_1(\cos\theta)\big)\,\right]_{-1}^{1},\\
\label{eq:ch13-dsigma-pred}
\diffXS^{\text{model}}(\theta) &=
u_{\text{norm}}\,\exp\!\big(1.10\,\text{inv\_T20}\,P_2(\cos\theta) +
0.60\,\text{inv\_T22}\,P_4(\cos\theta)\big).
\end{align}
\texttt{clamp} 在 \(\iTeleven\) 上是必要的：原则上 \(|\iTeleven| \le 1\)，
但 reduced-U 投影是「松弛公式」，会偶尔越界 \(\sim 10^{-2}\) 量级。系
数 \(\{1.20, -1.20, -0.20, 1.10, 0.60\}\) 是早期对 \(190\,\text{MeV}\)
基准做\textbf{形状拟合}（不是物理推导）得到的；它们在物理意义上对应
「\(P_2, P_4\) 角项的 phenomenological 系数」，未来应换为通过
\(C_{\ell m\ell' m'}^{J}\) 严格写下的 Wigner-D 项。

\lstinputlisting[style=python,
  caption={\texttt{compare\_Ay\_experiment.py} 的
           \texttt{\_predict\_observables\_from\_u}：reduced-U 不变量到
           \(\iTeleven, \diffXS\) 角度曲线},
  label={lst:ch13-predict-observables}]
  {code_excerpts/snippets/ch13_predict_observables.py}

\paragraph{Step G —— 主流程：选最近能量、生成报告。}
\texttt{main} 把上述步骤串成单次执行，并按
\(\big|\Tlab^{\text{best}} - \Tlab^{\text{target}}\big|\) 选最近能量、计
算误差度量、写 CSV/SVG/JSON/TXT 四类输出。

\lstinputlisting[style=python,
  caption={\texttt{compare\_Ay\_experiment.py} 的 \texttt{main} 主调
           （核心段）：选 \texttt{best\_energy} + 误差度量 + 写报告},
  label={lst:ch13-main-orchestration}]
  {code_excerpts/snippets/ch13_main_orchestration.py}

\subsection{\texttt{plot\_validation\_curves.py}：双面板对照图}
\label{subsec:ch13-code-plot}

\texttt{plot\_validation\_curves.py} 消费 \texttt{compare\_Ay\_experiment.py}
留下的 CSV/JSON 工件，用 matplotlib 画\textbf{双面板对照图}：左面板
\(\iTeleven\)（线性 y），右面板 \(\diffXS\)（log y）。每面板内嵌
\texttt{RMSE / MAE / max|err|} 文本注释，用于判别通过线（见
\cref{tab:ch13-thresholds}）。

\lstinputlisting[style=python,
  caption={\texttt{plot\_validation\_curves.py} 的双面板合成函数},
  label={lst:ch13-combined-plot}]
  {code_excerpts/snippets/ch13_combined_plot.py}

\paragraph{对照 origin 仓库。}
\texttt{tictac-origin/} 仓库里没有 \texttt{examples/} 目录，所有 Faddeev
驱动 + 比较都是 C++ 直出（\texttt{CPP/run} 直接吐 \texttt{output.txt}，
后期需要外部工具如 gnuplot 画图）。本章这套 Python 工作流是 Tic-tac
重构期间引入的、面向 RIKEN 实验的「快速迭代」轨道：研究者改一行
\(c_D, c_E\) 后跑 \texttt{compare\_Ay\_experiment.py --regenerate} 就可以
在 5\,min 内拿到 RMSE 是否改善的反馈，origin 路径需要一两个小时（手写
脚本 + 人眼对比）。

% ============================================================
\section{数字闭环}
\label{sec:ch13-numerical}

\subsection{\(190\,\text{MeV}\) 基准复现命令}
\label{subsec:ch13-reproduce-190}

最小复现命令（仓库根目录运行）：
\begin{lstlisting}[style=config, caption={190 MeV 极化观测量基准的复现脚本},
  label={lst:ch13-reproduce-cmd}]
# Step 1. 求解器（约 5 分钟，4 线程，Np_WP=Nq_WP=20）
python3 examples/deuteron_proton_Ay.py \
    --work-dir output/deuteron_proton_Ay \
    --target-tlab-mev 190 \
    --reuse-p123

# Step 2. 合成 + 比较实验（毫秒级）
python3 examples/compare_Ay_experiment.py \
    --work-dir output/deuteron_proton_Ay \
    --solver-out-dir output/deuteron_proton_Ay/solver_out \
    --target-tlab-mev 190 \
    --dsigma-unit fm2/sr

# Step 3. 画图（需要 matplotlib；anaroot-env 已包含）
micromamba run -n anaroot-env python3 examples/plot_validation_curves.py \
    --work-dir output/deuteron_proton_Ay
\end{lstlisting}

\subsection{在 \(\theta_{\text{cm}} = 30°/60°/90°/120°/150°\) 附近的数值表}
\label{subsec:ch13-numerical-table}

\paragraph{角度选择。}
实验文件 \texttt{T.txt}（Sekiguchi et al.\,2002，\(\Tlab = 190\,\text{MeV}/u\)
氘核作为入射极化粒子，\(d+p\) 弹性）覆盖
\(\theta_{\text{cm}} \in [38.9°, 164.8°]\)，没有 \(30°\) 数据点；本节挑选
五个最接近的角度：\(38.9°, 58.9°, 88.3°, 119.0°, 148.6°\)。表中模型列
取自 \texttt{output/deuteron\_proton\_Ay/best\_energy\_iT11\_curve.csv}
（求解器实际能量
\(\Tlab^{\text{best}} = 161.15\,\text{MeV}\)，能量偏离
\(\Delta\Tlab = 28.85\,\text{MeV}\)，落在 \(\pm 40\,\text{MeV}\) 通过线
内）。

\begin{numericalbox}
\textbf{Tic-tac 190 MeV 基准的 \(\iTeleven\) 数值闭环}\\[4pt]
（求解器配置：\(\NpWP=\NqWP=20\)、\texttt{LO\_internal} 势、
\texttt{tensor\_force=true}、\texttt{isospin\_breaking\_1S0=true}、
\texttt{two\_J\_3N\_max=1}、\texttt{J\_2N\_max=2}、4 线程；实际能量
\(\Tlab^{\text{best}} = 161.15\,\text{MeV}\)；
\(\iTeleven\) RMSE = 0.0959，MAE = 0.0789，max|err| = 0.184）
\begin{center}
\begin{tabular}{rccc}
\toprule
\(\theta_{\text{cm}}\,[\text{deg}]\)
& \(\iTeleven^{(\text{exp})} \pm \Delta\)
& \(\iTeleven^{(\text{model})}\) (本章实现)
& 残差 \(\Delta = \text{model} - \text{exp}\) \\
\midrule
38.9  & \(+0.304 \pm 0.004\) & \(+0.1767\) & \(-0.1273\) \\
58.9  & \(-0.061 \pm 0.004\) & \(-0.0061\) & \(+0.0549\) \\
88.3  & \(-0.289 \pm 0.006\) & \(-0.2581\) & \(+0.0309\) \\
119.0 & \(-0.238 \pm 0.013\) & \(-0.1044\) & \(+0.1336\) \\
148.6 & \(+0.012 \pm 0.006\) & \(+0.1227\) & \(+0.1107\) \\
\bottomrule
\end{tabular}
\end{center}

\textbf{读法。} 模型曲线在 \(\theta < 100°\) 区域形状对、幅度低
（\(\sim 60\%\) 实验幅度），与典型「\(A_y\) puzzle」一致；
\(\theta > 120°\) 区域出现号反向，说明只到 \(j_{2N}^{\max} = 2\)、
\(J_{3N}^{\max} = 1/2\) 与 reduced-U 投影都还不够支持后向角细节。
完整的 \(j_{2N}^{\max} = 5,
J_{3N}^{\max} = 7/2\) 计算（\cref{app:E_input_dictionary} 附录 E）外加 3NF 开关将把
RMSE 从 \(0.10\) 量级压到 \(0.02\) 量级。
\end{numericalbox}

\paragraph{\(\diffXS\) 数值表。}
取最接近 \(30°/60°/90°/120°/150°\) 的实验角度：
\(31.27°, 62.54°, 91.00°, 121.07°, 151.54°\)（来自
\texttt{DSigamaOverDOmega.txt} 第二列原文，单位
\(\text{mb}/\text{sr}\)；本表换算为 \(\text{fm}^2/\text{sr}\)
即除以 10）。

\begin{numericalbox}
\textbf{Tic-tac 190 MeV 基准的 \(\diffXS\) 数值闭环}\\[4pt]
（同求解器配置；\(\diffXS\) MAE = 2.12 \(\text{fm}^2/\text{sr}\)，
RMSE = 2.44，相对 RMSE = 83.87；模型与实验差距源自当前 reduced-U
\texttt{exp(...)} 形状拟合在中后向角的发散——见 \cref{subsec:ch13-pitfall-reducedU-vs-full}）
\begin{center}
\begin{tabular}{rcc}
\toprule
\(\theta_{\text{cm}}\,[\text{deg}]\)
& \(\diffXS^{(\text{exp})}\,[\text{fm}^2/\text{sr}]\)
& \(\diffXS^{(\text{model})}\,[\text{fm}^2/\text{sr}]\) (本章实现) \\
\midrule
31.27  & 0.4423  & 3.217  \\
62.54  & 0.0632  & 1.394  \\
91.00  & 0.0180  & 1.185  \\
121.07 & 0.0126  & 1.477  \\
151.54 & 0.0267  & 3.540  \\
\bottomrule
\end{tabular}
\end{center}

\textbf{读法。} 实验 \(\diffXS\) 量级跨 2 个数量级（前向 \(0.44\)、
中后向 \(0.013\)），而 reduced-U \texttt{exp(P2 + P4)} 拟合的预言量级
基本均匀（\(\sim 1\,\text{fm}^2/\text{sr}\)）——说明这个简化映射只对
极化\textbf{形状}有效，对截面\textbf{绝对值}还需要替换为完整的
Wigner-D 求迹（\cref{subsec:ch13-pitfall-reducedU-vs-full}）。这正是
为什么当前 \texttt{solver\_validation\_*.json} 里 \texttt{status} 标
为 \texttt{fail}（\texttt{dsigma\_rel\_rmse} 超过 \(0.05\) 通过线）：早
期开发期\textbf{有意保留}此 fail 状态作为「待重构」标记。
\end{numericalbox}

\subsection{对比图 PDF}
\label{subsec:ch13-numerical-figures}

\texttt{plot\_validation\_curves.py} 在
\texttt{output/deuteron\_proton\_Ay/} 下产出
\texttt{best\_energy\_iT11\_comparison.png} 与
\texttt{best\_energy\_dsigma\_comparison.png}（外加 SVG）。这两张
PNG 内嵌 RMSE/MAE/max|err| 文字注释，是日常调试用的「最小看图反
馈」。

\begin{figure}[h]
\centering
\fbox{\parbox{0.92\textwidth}{
\centering
\textbf{\textit{占位说明}}：图 \texttt{best\_energy\_iT11\_comparison.png}
（左）与 \texttt{best\_energy\_dsigma\_comparison.png}（右）。

实际嵌入路径：\\
\texttt{output/deuteron\_proton\_Ay/best\_energy\_iT11\_comparison.png}\\
\texttt{output/deuteron\_proton\_Ay/best\_energy\_dsigma\_comparison.png}

\medskip

\textit{典型形态}：左图显示模型曲线（红线）在 \(\theta \in
[40°, 120°]\) 区间与实验数据点（黑点）形状基本同步、幅度低 \(\sim
40\%\)；右图模型曲线（红线，log y）远离实验点（黑点）2 个量级，
体现 reduced-U 投影对绝对截面的失真。
}}
\caption{\(\Tlab^{\text{target}} = 190\,\text{MeV}\)、\(\Tlab^{\text{best}}
= 161.15\,\text{MeV}\) 的 \(\iTeleven\)、\(\diffXS\) 曲线对比（reduced-U
模型 vs.\ Sekiguchi et al.\,2002）。}
\label{fig:ch13-comparison}
\end{figure}

\subsection{通过线（pass thresholds）}
\label{subsec:ch13-thresholds}

\begin{table}[h]
\centering
\caption{\texttt{compare\_Ay\_experiment.py} 默认通过线（来自
\texttt{--ay-rmse-pass}、\texttt{--dsigma-rel-rmse-pass}、
\texttt{--energy-delta-pass}）。}
\label{tab:ch13-thresholds}
\begin{tabular}{lccc}
\toprule
观测量 / 维度 & 通过线 & 警告线 & 当前 \(190\,\text{MeV}\) 实测 \\
\midrule
\(\iTeleven\) RMSE          & \(\le 0.02\) & \(\le 0.05\) & \(0.0959\) (fail) \\
\(\diffXS\) 相对 RMSE       & \(\le 0.05\) & \(\le 0.10\) & \(83.87\) (fail) \\
\(|\Delta\Tlab|\) (MeV)     & \(\le 40\)   & \(\le 80\)   & \(28.85\) (pass) \\
\bottomrule
\end{tabular}
\end{table}

\paragraph{为什么当前 status = fail 但仍然「合格」。}
仓库的开发约定是：\textbf{在重构窗口内允许 fail，但不允许「无诊断信息
的 fail」}。\texttt{solver\_validation\_*.json} 中的 \texttt{best\_energy}
块、reduced-U 不变量、\texttt{phase\_sign} 都被持续记录，作为下一阶段
（用完整 Wigner-D 求迹替换 reduced-U 投影）的基线对照。
\cref{ch:14_miller_gate1} 之后的 J-tensor 分析能力实现一旦上线，预期
\texttt{status} 会自动转 pass。

% ============================================================
\section{接口与依赖}
\label{sec:ch13-interface}

\subsection{上游：\texttt{U\_PW\_elements\_*.txt} 文件格式}
\label{subsec:ch13-upstream}

求解器输出文件名形如
\verb|U_PW_elements_Np_<NpWP>_Nq_<NqWP>_JP_<2J>_<P>_Jmax_<j2N_max>_PSI_<idx>.txt|，
对每行存
\begin{equation}
\Tlab,\; E_{\text{cm}},\; q_{\text{idx}},\;
U_{00},\; U_{01},\; U_{10},\; U_{11},
\end{equation}
其中 \(U_{ij}\) 是用 Python \texttt{complex(...)} 直接可解析的复数字符串
（如 \texttt{(1.234,-5.678)}）。\texttt{parse\_u\_file} 是这一格式的
\textbf{唯一定义}——若求解器升级输出列数，更新这里即可。

\subsection{下游：四类工件文件}
\label{subsec:ch13-downstream}

\begin{itemize}
  \item \texttt{best\_energy\_iT11\_curve.csv}、
        \texttt{best\_energy\_dsigma\_curve.csv} —
        模型 vs.\ 实验的角度曲线，被
        \texttt{plot\_validation\_curves.py} 与
        \texttt{plot\_3nf\_effect.py} 共用。
  \item \texttt{best\_energy\_iT11\_comparison.svg}、
        \texttt{best\_energy\_dsigma\_comparison.svg} —
        独立 SVG 图（不依赖 matplotlib，纯 Python 字符串拼接）。
  \item \texttt{solver\_validation\_tlab\_<label>.json} —
        机器可读全摘要，被回归测试
        \texttt{tests/test\_190mev\_data\_pipeline.py} 解析。
  \item \texttt{solver\_validation\_tlab\_<label>.txt} —
        人类可读摘要（CI 日志友好）。
\end{itemize}

\subsection{第 12 章共用振幅。第 14 章共用 \(\Uop\)-矩阵元。}
\texttt{compare\_Ay\_experiment.py} 与
\cref{ch:12_observables_basics} 的 \texttt{deuteron\_proton\_dsigma.py}
（如有）共享同一份 \texttt{U\_PW\_elements\_*.txt}；
\cref{ch:14_miller_gate1} 的 Miller-gate 1 相移工具也读同一份文件。
合理的工程边界是：\textbf{合成代码彼此独立、可任意替换}（如
reduced-U → 完整 Wigner-D），但 \(\Uop\)-文件格式被冻结，由
\texttt{src/io/write\_pw\_amplitude.cpp}（C++ 端）与
\texttt{examples/solver\_u\_file\_utils.py}（Python 端）联合定义。

\subsection{对照表：origin vs.\ Tic-tac}
\label{subsec:ch13-origin-comparison}

\begin{table}[h]
\centering
\caption{\(\Uop\)-矩阵元到极化曲线的工作流：origin vs.\ 当前 Tic-tac。}
\label{tab:ch13-origin-vs-current}
\begin{tabular}{lll}
\toprule
阶段 & \texttt{tictac-origin} & 当前 \texttt{Tic-tac} \\
\midrule
求解器调度 & 手写 shell 脚本 & \texttt{deuteron\_proton\_Ay.py} \\
合成 & \(\diffXS\) 由 C++ 直出 & \texttt{compare\_Ay\_experiment.py} \\
\(\iTeleven\) 计算 & 不在 origin 中 & 本章实现（reduced-U） \\
张量观测量 & 不在 origin 中 & 本章 \(\Ttwentyzero, \ldots\) inv 量 \\
对实验比较 & gnuplot 手画 & 自动 CSV/SVG/JSON/TXT \\
3NF 开关扫描 & 重新编译 & \texttt{run\_3nf\_sweep.py} 一行 \\
回归测试 & 无 & \texttt{tests/test\_190mev\_data\_pipeline.py} \\
\bottomrule
\end{tabular}
\end{table}

% ============================================================
\section{常见陷阱}
\label{sec:ch13-pitfalls}

\subsection{陷阱 A：\(\Ay\) 与 \(\iTeleven\) 的 \(\sqrt{3/2}\) 因子}
\label{subsec:ch13-pitfall-Ay-vs-iT11}

\begin{pitfallbox}
\textbf{症状}：你拿 Sekiguchi et al.\,2002 表 III 的数据列对比
\texttt{quick\_Ay\_test.py} 输出，发现实验值整体小 \(\sim 22\%\)，
似乎模型「准确度过高」。

\textbf{诊断}：检查源数据列头是 \(\iTeleven\) 还是 \(A_y\)。
Sekiguchi et al.\,的 \texttt{T.txt} 列头是 \(\iTeleven\)（氘核作为入
射极化粒子时的矢量分析能力）；旧 RIKEN 数据有时直接列 \(A_y^{(d)}\)。

\textbf{修复}：式 \eqref{eq:ch13-Ay-iT11}：
\(A_y^{(d)} = \sqrt{3/2}\,\iTeleven \approx 1.2247\,\iTeleven\)。代码
中如果想转换，乘 \(\sqrt{1.5}\)。注意\textbf{不是} \(\sqrt{2}\)（这是
另一种与 \(s_a = 1/2\) 有关的归一）。

\textbf{为什么是 \(\sqrt{3/2}\)}：从
\(\tau_{1,+1}^{(s_a=1)} = -\frac{1}{\sqrt{2}}(s_x + i s_y)\)
（\cref{eq:ch13-Tkq-master}），加上 \(\sqrt{2k+1} = \sqrt{3}\) 的张量
归一，\(\Ay = \langle s_y\rangle\) 直接对应；除一下得
\(\sqrt{3/2}\)。
\end{pitfallbox}

\subsection{陷阱 B：Madison 约定与号倒错}
\label{subsec:ch13-pitfall-madison}

\begin{pitfallbox}
\textbf{症状}：模型 \(\iTeleven\) 与实验 \(\iTeleven\) 在所有角度同号
反向（前向正变成前向负），但绝对值还能基本对上。

\textbf{诊断}：检查 \(\hat{y}\) 的取向。Madison 约定：
\(\hat{y} = \hat{k}_{\text{in}} \times \hat{k}_{\text{out}} / \|\cdots\|\)。
有些老论文（特别是 Bonn 70 年代实验）用反向：
\(\hat{y} = \hat{k}_{\text{out}} \times \hat{k}_{\text{in}} / \|\cdots\|\)。

\textbf{修复}：在 \texttt{compare\_Ay\_experiment.py} 里
\texttt{phase\_sign} 字段就是为了诊断这件事保留的。如发现
\texttt{phase\_sign = -1.0}，说明模型 \(\iTeleven\) 与实验差一个全局号
（\(\Ttwentyone\) 同样翻号；\(\Ttwentyzero, \Ttwentytwo\) 不变）。修
复办法：在 \cref{eq:ch13-iT11-pred} 上乘 \texttt{phase\_sign}（或更彻
底地，重定义 \(\tau_{1,+1}\) 的相位）。
\end{pitfallbox}

\subsection{陷阱 C：reduced-U 投影 vs.\ 完整 Wigner-D 求迹}
\label{subsec:ch13-pitfall-reducedU-vs-full}

\begin{pitfallbox}
\textbf{症状}：\(\iTeleven\) RMSE 还能压在 \(0.10\) 量级，但
\(\diffXS\) 模型曲线和实验曲线差 \(2\sim 3\) 数量级（如
\cref{subsec:ch13-numerical-table} 所示，模型 \(\sim 1\,\text{fm}^2/\text{sr}\)、
实验 \(\sim 0.018\,\text{fm}^2/\text{sr}\) 在 \(91°\) 处）。

\textbf{诊断}：当前 \texttt{\_predict\_observables\_from\_u} 在
\(\diffXS\) 上用的是 \texttt{u\_norm * exp(1.10 inv\_T20 P\_2 + 0.60
inv\_T22 P\_4)} 的「形状拟合」，没有把 \(M\)-矩阵的
spin-helicity 求迹结构正确折进来；\(u_{\text{norm}}\) 也没有乘
\(\big|f(p, q)\big|^{-2}\) 的相空间因子。

\textbf{修复路径}：
\begin{enumerate}
  \item 把 \texttt{\_predict\_observables\_from\_u} 重写成「Wigner-D 求
        和 + spin tensor 求迹」的完整教科书形式
        （\cref{eq:ch13-Tkq-master}）；
  \item \(u_{\text{norm}}\) 乘上正确的运动学因子
        \((q_{\text{out}} / q_{\text{in}}) \cdot |\mu|^2\)（参见
        \cref{ch:12_observables_basics}）；
  \item 完整通道集合 \((J^\pi = \tfrac{1}{2}^{\pm}, \tfrac{3}{2}^{\pm},
        \tfrac{5}{2}^{\pm}, \ldots)\) 全部参与求和，而不是只
        \((J^\pi = 1^\pm)\)。
\end{enumerate}

\textbf{当前合理化}：在 reduced-U 实现下保留
\texttt{status = fail} 是\emph{有意为之}（\cref{subsec:ch13-thresholds}），
等 J-tensor 实现切换上线 \texttt{status} 自然变 pass。
\end{pitfallbox}

\subsection{陷阱 D：输入 \(\Tlab\) ≠ 求解器实际 \(\Tlab\)}
\label{subsec:ch13-pitfall-input-vs-solver-tlab}

\begin{pitfallbox}
\textbf{症状}：你 \texttt{--target-tlab-mev 190} 跑完后看
\texttt{solver\_validation\_tlab\_190MeV.json}，发现
\texttt{best\_tlab\_mev = 161.145}，差了将近
\(\Delta\Tlab \approx 29\,\text{MeV}\)。

\textbf{诊断}：Tic-tac 求解器内部把 \(\Tlab\) 「投影」到\textbf{ on-shell
q-bin 的中点}。原始请求的 \(190\,\text{MeV}\) 对应一个连续的 \(q\) 值；
但 WP 离散把 \([0, q_{\max}]\) 切成 \(\NqWP\) 个 bin，求解器只能在 bin 中
点上算 \(\Uop\)-矩阵元。Tic-tac 的当前选择是「取最近 bin 中点」，所以
真正算出来的能量可能与请求差 \(\pm 10\sim 40\,\text{MeV}\)。

\textbf{修复}：
\begin{enumerate}
  \item 把 \(\NqWP\) 加大到 \(40\) 或 \(80\)，bin 间距压到
        \(\sim 5\,\text{MeV}\) 量级；
  \item 检查 \texttt{src/config/run\_organizer.cpp::find\_on\_shell\_bins}
        看哪个 bin 被选中；
  \item 如果一定要精确 \(190\,\text{MeV}\)，可以用 \texttt{--chebyshev-s}
        与 \texttt{--chebyshev-t} 调整网格分布让某个 bin 中点正好落在
        \(190\,\text{MeV}\)。
\end{enumerate}

\textbf{为什么默认行为是「最近 bin」}：因为 Faddeev 求解器的设计是
\textbf{批量算多能量}，逐能量做 root-finding 会把成本翻好几倍。在重构
窗口内、误差 \(\le 40\,\text{MeV}\) 是\textbf{通过线}（\cref{tab:ch13-thresholds}），
反映了「\(\iTeleven\) 在
\([150, 230]\,\text{MeV}\) 区间形状不剧烈变化」这个物理事实。
\end{pitfallbox}

\subsection{陷阱 E：\(T_{kq}\) 的 Hermitian 性与实/虚部约定}
\label{subsec:ch13-pitfall-hermiticity}

\begin{pitfallbox}
\textbf{症状}：你想从 \(\tau_{2,+2}\) 算 \(\Ttwentytwo\)，但发现取 \(M\)-
矩阵的「实部」时和 Sekiguchi 表里的 \(\Ttwentytwo\) 差一个相位。

\textbf{诊断}：\(\tau_{kq}\) 不是 Hermitian，它满足
\(\tau_{kq}^\dagger = (-1)^{q}\,\tau_{k,-q}\)。所以 \(\Tr[M\tau_{kq}M^\dagger]\)
本身\textbf{是复数}；只有「\(\tau_{k,q} + (-1)^{q}\tau_{k,-q}\)」与
「\(\tau_{k,q} - (-1)^{q}\tau_{k,-q}\)」的两个 Hermitian 组合给出实数
观测量。

\textbf{修复}：实验上的 \(\Re\Ttwentytwo\) 对应
\(\Tr[M(\tau_{2,+2} + \tau_{2,-2})M^\dagger]/(2\Tr[MM^\dagger])\)，
\(\Im\Ttwentytwo\) 对应另一种组合。Madison 约定下当散射平面取
\(xz\) 平面（\(\phi=0\)）时，\(\Im\Ttwentytwo = 0\)，所以实验报的就是
\(\Re\Ttwentytwo\)。代码里 \texttt{inv\_T22} 取的正是
\(\Re[U_{00}^{*}U_{11} - U_{01}^{*}U_{10}]\)（实部），与这个约定一致。
\end{pitfallbox}

\subsection{陷阱 F：\(\Ttwentyone\) 在 \(\theta = 0\) 处必须为零，
但 reduced-U 投影会越界}
\label{subsec:ch13-pitfall-T21-at-zero}

\begin{pitfallbox}
\textbf{症状}：\(\Ttwentyone(\theta = 0)\) 物理上必为零（散射平面在
正向角无定义，\(T_{2,\pm 1}\) 受限于 \(\sin\theta\) 因子），但 reduced-U
投影输出可能有非零值。

\textbf{诊断}：reduced-U 投影里 \(\Ttwentyone\) 没有显式 \(\sin\theta\)
因子（仅有 \(P_2 \cos^2\theta\) 项），导致 \(\theta = 0\) 时未受约束。

\textbf{修复}：在 \cref{eq:ch13-iT11-pred} 后添加
\(\Ttwentyone^{\text{model}}(\theta) = \text{const} \cdot \text{inv\_T21}
\cdot \sin(2\theta)\)（这是 reduced-U 的「显式 \(\sin\) 修正」）。当前
\texttt{compare\_Ay\_experiment.py} 仅算 \(\iTeleven\) 与
\(\diffXS\)；\(\Ttwentyzero, \Ttwentyone, \Ttwentytwo\) 暂时只输出
不变量到 JSON、不画曲线。等完整 J-tensor 实现上线后这个 \(\sin\theta\)
因子会从 Wigner-D 求迹自然涌现。
\end{pitfallbox}

\subsection{陷阱 G：单位与 \(\diffXS\) 的 mb/sr ↔ fm²/sr}
\label{subsec:ch13-pitfall-units}

\begin{pitfallbox}
\textbf{症状}：你画 \(\diffXS\) 看到模型曲线在 \(\sim 1\) 量级、实
验曲线在 \(\sim 0.018\) 量级，相差 \(\sim 50\) 倍。

\textbf{诊断}：检查单位。\texttt{1 mb/sr = 0.1 fm²/sr}（\(1\,\text{mb}
= 0.1\,\text{fm}^2\)）。\texttt{compare\_Ay\_experiment.py} 默认
输出 \texttt{mb/sr}，实验文件
\texttt{DSigamaOverDOmega.txt} 的源单位也是 \texttt{mb/sr}（文件头自
带说明），所以默认配置下不会差 \(50\) 倍。但如果你用
\texttt{--dsigma-unit fm2/sr} 而模型仍按 \texttt{mb/sr} 公式算，
就会差 \(10\) 倍——这是\texttt{convert\_dsigma\_value} 设计存在的
原因。

\textbf{修复}：永远显式传 \texttt{--dsigma-unit}。CI 在
\texttt{tests/test\_190mev\_data\_pipeline.py} 里强制断言
\texttt{units.dsigma\_output == "mb/sr"}（默认）以防漂移。
\end{pitfallbox}

\subsection{陷阱 H：\(P_{123}\) 缓存与 \texttt{--reuse-p123}}
\label{subsec:ch13-pitfall-p123-cache}

\begin{pitfallbox}
\textbf{症状}：你改了 \texttt{potential\_model} 重跑
\texttt{deuteron\_proton\_Ay.py --reuse-p123}，发现 \(\iTeleven\) 几乎不
变。

\textbf{诊断}：\texttt{--reuse-p123} 让求解器跳过 \(P_{123}\) 重新计算
（\(P_{123}\) 矩阵本身只依赖网格与角动量截断，不依赖势模型）；但势矩阵
\(V\) 与 \(\Uop\)-矩阵元仍重新算。如果 \(\iTeleven\) 几乎不变，可能是：
\begin{enumerate}
  \item 势模型 \emph{真的} 在 \(190\,\text{MeV}\) 上对 \(\iTeleven\) 不
        敏感（这种情况要看哪些通道占主）；
  \item 实际你忘了去掉 \texttt{--reuse-p123} 同时也没改其他参数；
  \item 求解器输出目录被其它进程覆盖。
\end{enumerate}

\textbf{修复}：每次实验性切换势模型时，\textbf{第一次}用
\texttt{calculate\_and\_store\_P123=true}（即不加 \texttt{--reuse-p123}）
让缓存重建，确认 \texttt{p123/} 目录的时间戳更新；之后才能 \texttt{--reuse-p123}。
\end{pitfallbox}

\section*{本章小结}

本章把第 \ref{ch:11_faddeev_solver} 章求出的
\(\Uop\)-矩阵元一路推到「与 Sekiguchi et al.\,2002 实验数据可比的角度
曲线」：

\begin{itemize}
  \item \cref{sec:ch13-motivation,sec:ch13-math} 给出从极化态密度矩阵 +
        球张量基出发到张量分析能力 \(\Ttwentyzero, \Ttwentyone,
        \Ttwentytwo\) 与矢量分析能力 \(\iTeleven\) 的求迹公式
        \eqref{eq:ch13-Tkq-master}，并强调与 \(\Ay^{(d)}\) 之间的因子
        \(\sqrt{3/2}\) （\cref{eq:ch13-Ay-iT11}）；
  \item \cref{sec:ch13-discretization,sec:ch13-algorithm} 把流水线
        拆成 Solver / 实验读取 / 合成三段，并给出
        \cref{alg:ch13-pol-obs} 的伪码骨架；
  \item \cref{sec:ch13-code} 走完
        \texttt{deuteron\_proton\_Ay.py} +
        \texttt{compare\_Ay\_experiment.py} +
        \texttt{plot\_validation\_curves.py} 三个 Python 脚本的核心函数；
  \item \cref{sec:ch13-numerical} 在 \(190\,\text{MeV}\) 基准下给出
        \(\iTeleven, \diffXS\) 在五个代表角度处的数值闭环表，并明确指出
        当前 reduced-U 实现在 \(\diffXS\) 绝对值上不通过线、但作为
        \(\iTeleven\) 形状的早期诊断有效；
  \item \cref{sec:ch13-pitfalls} 列了 8 个具体陷阱（\(\Ay\) vs
        \(\iTeleven\) 因子、Madison 约定号、reduced-U vs.\ 完整 Wigner-D、
        \(\Tlab\) 投影偏差、张量算符 Hermitian 性、\(T_{21}\) 在
        \(\theta=0\) 必为零、单位换算、\(P_{123}\) 缓存）。
\end{itemize}

下一章 \ref{ch:14_miller_gate1} 把同一份 \(\Uop\)-矩阵元用作
Miller-gate 1 相移提取，对照 nd 散射相移基准；
\cref{ch:15_3nf_physics} 起则在本章建立的极化通道上系统地切换
\(c_D, c_E, c_1, c_3, c_4\)，量化 3NF 对 \(\iTeleven\) puzzle 的修正。
